% GENERATED FILE. Edit canonical lessons, then run npm run generate:books.
% canonical-source-hash: fnv1a64:d4b06fa3aa1e4022
% canonical-lessons: UR-C408-simt, UR-C408-shimal, UR-C408-junub, UR-C408-mashriq, UR-C408-maghrib

\chapter{North, South, East and West}
\label{ch:ur-north-south-east-and-west}
\hlchaptermodality{408}

\begin{chapteropening}
\textbf{By the end of this chapter:} \emph{I can say which direction a place lies in.}

Complete the last lesson of chapter 408: I can say which direction a place lies in.
\end{chapteropening}

\section[\texorpdfstring{\ur{سمت} (simt)}{سمت (simt)}]{\ur{سمت} (simt) — a direction}

\label{lesson:UR-C408-simt}



\begin{quote}
Before the new one: say the Urdu for a repair, then the Urdu for to move (something) aside.
\end{quote}

\subsection*{You'll want to know: \ur{سمت}}
\textbf{\ur{سمت}} — \emph{simt} — \textquotedblleft{}a direction\textquotedblright{}.

\begin{quote}
\ur{اسٹیشن} \ur{کس} \ur{سمت} \ur{میں} \ur{ہے؟}
\end{quote}

\begin{quote}
\emph{isṭeshan kis simt meṅ hai?}
\end{quote}

\begin{quote}
Which direction is the station in?
\end{quote}

Feminine.

\begin{cousinweb}
\textbf{\ur{سمت}} (\emph{simt}) is usually traced, through Persian, to Arabic \emph{samt}, \textquotedblleft{}a way, a direction\textquotedblright{}.
\end{cousinweb}

\begin{grammarlens}[title={What you've built}]
One more word for getting around.
\end{grammarlens}

\subsection*{Your turn}
\begin{itemize}
\raggedright
  \item \emph{Say it:} \emph{simt}
  \item \emph{Say it:} \emph{simt}, once more
  \item \emph{Say it:} \emph{isṭeshan kis simt meṅ hai?} — Which direction is the station in?
  \item \emph{Recall:} say the Urdu for a repair, then the Urdu for to move (something) aside, then say \emph{simt} again
\end{itemize}

\begin{tcolorbox}[breakable,colback=teal!4,colframe=teal!35!black,title={Before you move on}]
What does \ur{سمت} mean? (\textquotedblleft{}A direction\textquotedblright{}.) Say it once more.
\end{tcolorbox}
\section[\texorpdfstring{\ur{شمال} (shimāl)}{شمال (shimāl)}]{\ur{شمال} (shimāl) — the north}

\label{lesson:UR-C408-shimal}



\begin{quote}
Before the new one: say the Urdu for to move (something) aside, then the Urdu for a direction.
\end{quote}

\subsection*{You'll want to know: \ur{شمال}}
\textbf{\ur{شمال}} — \emph{shimāl} — \textquotedblleft{}the north\textquotedblright{}.

\begin{quote}
\ur{پہاڑ} \ur{شہر} \ur{کے} \ur{شمال} \ur{میں} \ur{ہیں۔}
\end{quote}

\begin{quote}
\emph{pahāṛ shahr ke shimāl meṅ haiṅ.}
\end{quote}

\begin{quote}
The mountains lie to the north of the city.
\end{quote}

Masculine.

\begin{cousinweb}
\textbf{\ur{شمال}} (\emph{shimāl}) is usually traced, through Persian, to Arabic \emph{shamāl}, \textquotedblleft{}the north\textquotedblright{}.
\end{cousinweb}

\begin{grammarlens}[title={What you've built}]
One more word for getting around.
\end{grammarlens}

\subsection*{Your turn}
\begin{itemize}
\raggedright
  \item \emph{Say it:} \emph{shimāl}
  \item \emph{Say it:} \emph{shimāl}, once more
  \item \emph{Say it:} \emph{pahāṛ shahr ke shimāl meṅ haiṅ.} — The mountains lie to the north of the city.
  \item \emph{Recall:} say the Urdu for to move (something) aside, then the Urdu for a direction, then say \emph{shimāl} again
\end{itemize}

\begin{tcolorbox}[breakable,colback=teal!4,colframe=teal!35!black,title={Before you move on}]
What does \ur{شمال} mean? (\textquotedblleft{}The north\textquotedblright{}.) Say it once more.
\end{tcolorbox}
\section[\texorpdfstring{\ur{جنوب} (junūb)}{جنوب (junūb)}]{\ur{جنوب} (junūb) — the south}

\label{lesson:UR-C408-junub}



\begin{quote}
Before the new one: say the Urdu for a direction, then the Urdu for the north.
\end{quote}

\subsection*{You'll want to know: \ur{جنوب}}
\textbf{\ur{جنوب}} — \emph{junūb} — \textquotedblleft{}the south\textquotedblright{}.

\begin{quote}
\ur{سمندر} \ur{جنوب} \ur{کی} \ur{سمت} \ur{میں} \ur{ہے۔}
\end{quote}

\begin{quote}
\emph{samandar junūb kī simt meṅ hai.}
\end{quote}

\begin{quote}
The sea lies to the south.
\end{quote}

Masculine.

\begin{cousinweb}
\textbf{\ur{جنوب}} (\emph{junūb}) is usually traced, through Persian, to Arabic \emph{janūb}, \textquotedblleft{}the south\textquotedblright{}.
\end{cousinweb}

\begin{grammarlens}[title={What you've built}]
One more word for getting around.
\end{grammarlens}

\subsection*{Your turn}
\begin{itemize}
\raggedright
  \item \emph{Say it:} \emph{junūb}
  \item \emph{Say it:} \emph{junūb}, once more
  \item \emph{Say it:} \emph{samandar junūb kī simt meṅ hai.} — The sea lies to the south.
  \item \emph{Recall:} say the Urdu for a direction, then the Urdu for the north, then say \emph{junūb} again
\end{itemize}

\begin{tcolorbox}[breakable,colback=teal!4,colframe=teal!35!black,title={Before you move on}]
What does \ur{جنوب} mean? (\textquotedblleft{}The south\textquotedblright{}.) Say it once more.
\end{tcolorbox}
\section[\texorpdfstring{\ur{مشرق} (mashriq)}{مشرق (mashriq)}]{\ur{مشرق} (mashriq) — the east}

\label{lesson:UR-C408-mashriq}



\begin{quote}
Before the new one: say the Urdu for the north, then the Urdu for the south.
\end{quote}

\subsection*{You'll want to know: \ur{مشرق}}
\textbf{\ur{مشرق}} — \emph{mashriq} — \textquotedblleft{}the east\textquotedblright{}.

\begin{quote}
\ur{سورج} \ur{مشرق} \ur{سے} \ur{نکلتا} \ur{ہے۔}
\end{quote}

\begin{quote}
\emph{sūraj mashriq se nikaltā hai.}
\end{quote}

\begin{quote}
The sun rises in the east.
\end{quote}

Masculine.

\begin{cousinweb}
\textbf{\ur{مشرق}} (\emph{mashriq}) is usually traced, through Persian, to Arabic \emph{mashriq}, \textquotedblleft{}the place where the sun rises\textquotedblright{}.
\end{cousinweb}

\begin{grammarlens}[title={What you've built}]
One more word for getting around.
\end{grammarlens}

\subsection*{Your turn}
\begin{itemize}
\raggedright
  \item \emph{Say it:} \emph{mashriq}
  \item \emph{Say it:} \emph{mashriq}, once more
  \item \emph{Say it:} \emph{sūraj mashriq se nikaltā hai.} — The sun rises in the east.
  \item \emph{Recall:} say the Urdu for the north, then the Urdu for the south, then say \emph{mashriq} again
\end{itemize}

\begin{tcolorbox}[breakable,colback=teal!4,colframe=teal!35!black,title={Before you move on}]
What does \ur{مشرق} mean? (\textquotedblleft{}The east\textquotedblright{}.) Say it once more.
\end{tcolorbox}
\section[\texorpdfstring{\ur{مغرب} (maghrib)}{مغرب (maghrib)}]{\ur{مغرب} (maghrib) — the west}

\label{lesson:UR-C408-maghrib}



\begin{quote}
Before the new one: say the Urdu for the south, then the Urdu for the east.
\end{quote}

\subsection*{You'll want to know: \ur{مغرب}}
\textbf{\ur{مغرب}} — \emph{maghrib} — \textquotedblleft{}the west\textquotedblright{}.

\begin{quote}
\ur{شام} \ur{کو} \ur{سورج} \ur{مغرب} \ur{میں} \ur{چھپتا} \ur{ہے۔}
\end{quote}

\begin{quote}
\emph{shām ko sūraj maghrib meṅ chhuptā hai.}
\end{quote}

\begin{quote}
In the evening the sun goes down in the west.
\end{quote}

Masculine.

\begin{cousinweb}
\textbf{\ur{مغرب}} (\emph{maghrib}) is usually traced, through Persian, to Arabic \emph{maghrib}, \textquotedblleft{}the place where the sun sets\textquotedblright{}; it is also the name of the prayer at sunset.
\end{cousinweb}

\begin{grammarlens}[title={What you've built}]
One more word for getting around.
\end{grammarlens}

\subsection*{Your turn}
\begin{itemize}
\raggedright
  \item \emph{Say it:} \emph{maghrib}
  \item \emph{Say it:} \emph{maghrib}, once more
  \item \emph{Say it:} \emph{shām ko sūraj maghrib meṅ chhuptā hai.} — In the evening the sun goes down in the west.
  \item \emph{Recall:} say the Urdu for the south, then the Urdu for the east, then say \emph{maghrib} again
\end{itemize}

\begin{tcolorbox}[breakable,colback=teal!4,colframe=teal!35!black,title={Before you move on}]
What does \ur{مغرب} mean? (\textquotedblleft{}The west\textquotedblright{}.) Say it once more.
\end{tcolorbox}
